\subsection{局部紧子空间上的诱导测度}

设 X 是一个局部紧空间，\(\mu\) 是 X 上的一个测度，Y 是 X 的一个局部紧子空间。由于 Y 是 X 中一个开集与一个闭集的交（GT, I, §9, No.~7, 命题 12），故它是 \(\mu\)-可测的（No.~1, 命题 3 的推论）。对每个函数 \(g \in \mathcal{K}(\mathrm{Y};\mathbf{C})\)，设 \(g'\) 是在整个 X 上定义的、在 Y 上等于 g 而在 X - Y 上等于 0 的函数；我们来证明 \(g'\) 是 \(\mu\)-可积的。我们可以只限于 g 为实且 \(\geqslant 0\) 的情形（把 g 写成这类函数的线性组合即可）；由于 \(g'\) 有界且具有紧支集，只需证明 \(g'\) 是 \(\mu\)-可测的（No.~6, 定理 5）；但这由 \(g'\) 在 X 上为上半连续这一事实得出（No.~5, 命题 8 的推论）。于是我们可以给出下面的定义：

定义 4. --- 设 Y 是局部紧空间 X 的一个局部紧子空间，我们把由 X 上的测度 \(\mu\) 在 Y 上诱导的、记作 \(\mu_{Y}\) 或 \(\mu|Y\) 的测度理解为由下列公式定义的测度

\[
\int g d \mu_ {\mathrm{Y}} = \int g ^ {\prime} d \mu\tag{1}
\]

此公式对每个函数 \(g \in \mathcal{K}(\mathrm{Y};\mathbf{C})\) 成立，其中 \(g'\) 表示等于 \(g\)（在 \(\mathrm{Y}\) 上）、等于 0（在 \(\mathrm{X} - \mathrm{Y}\) 上）的函数。

例. --- 设 \(\mu\) 是 R 上的 Lebesgue 测度，I 是 R 中的任一区间；I 是 R 的一个局部紧子空间，且由 \(\mu\) 在 I 上诱导的测度是线性形式

\[
g \mapsto \int_ {a} ^ {b} g (x) d x
\]

（即 \(\mathcal{K}(\mathrm{I};\mathbf{C})\) 上的线性形式），其中 \(a\) 与 \(b\) 是 I 的端点（有限或无限）（参见 § 4, No.~4 的例）；换句话说，它就是我们已称之为 I 上的 Lebesgue 测度的那个测度。

当 Y 是 X 的开子空间时，定义 4 与 Ch. III, §2, No.~1 中给出的由 \(\mu\) 在 Y 上诱导的测度（或 \(\mu\) 在 Y 上的限制）的定义一致：事实上，对每个函数 \(g \in \mathcal{K}(\mathrm{Y}; \mathbf{C})\)，函数 \(g'\) 这时在 X 上连续。

关于诱导测度的积分，我们将在 Ch. V, §7 中详细研究，在此之前我们只需要下列结果：

引理 2. --- 设 \(\mu\) 是 \(X\) 上的一个正测度，\(K\) 是 \(X\) 的一个紧子集。

\begin{enumerate}
\def\labelenumi{(\roman{enumi})}
\item
  对 K 的每个紧（相应地，开）子集 H，\(\mu_{\mathrm{K}}(\mathrm{H}) = \mu (\mathrm{H})\)
\item
  为使 \(\mathbf{N}\)（\(\mathbf{K}\) 的子集）是 \(\mu_{\mathbf{K}}\)-可忽略的，其充要条件是它是 \(\mu\)-可忽略的。
\item
  若 S 是 \(\mu_{\mathrm{K}}\) 的支集，则 \(\operatorname{Supp}(\mu_{\mathrm{S}}) = \mathrm{S}\) 。
\setcounter{enumi}{0}
\item
  我们可以只限于 H 为紧集的情形。用 f 表示 H 在空间 K 中的特征函数；f 是上半连续的，因而是递减有向族 \((g_{\alpha})\)（由 \(\mathcal{K}_{+}(\mathrm{K})\) 中的函数组成）的下包络；我们有 \(\mu_{\mathrm{K}}(\mathrm{H}) = \inf_{\alpha} \int g_{\alpha} d\mu_{\mathrm{K}}\)（§4, No.~4, 命题 5 的推论 2）。若 \(g_{\alpha}^{\prime}\) 是在 K 上等于 \(g_{\alpha}\) 而在 X - K 上等于 0 的函数，则 \(g_{\alpha}^{\prime}\) 上半连续，且递减有向族 \((g_{\alpha}^{\prime})\) 的下包络是 H 在空间 X 中的特征函数 \(\varphi_{H}\)；因此
\end{enumerate}

\[
\mu (\mathrm{H}) = \inf _ {\alpha} \int g _ {\alpha} ^ {\prime} d \mu = \inf _ {\alpha} \int g _ {\alpha} d \mu_ {\mathrm{K}} = \mu_ {\mathrm{K}} (\mathrm{H})
\]

由 (1) 即得。

\begin{enumerate}
\def\labelenumi{(\roman{enumi})}
\setcounter{enumi}{1}
\item
  若 N 是 \(\mu\)-可忽略的，则对每个 \(\varepsilon > 0\)，存在 N 在 X 中的一个相对紧的开邻域 U，使得 \(\mu(\mathrm{U}) \leqslant \varepsilon\)；由于 \(\mathrm{K} - (\mathrm{U} \cap \mathrm{K})\) 是紧的，由 (i) 得 \(\mu_{\mathrm{K}}(\mathrm{U} \cap \mathrm{K}) \leqslant \mu(\mathrm{U}) \leqslant \varepsilon\)，从而 N 是 \(\mu_{K}\)-可忽略的。反之，若 N 是 \(\mu_{K}\)-可忽略的，则存在 N 在 X 中的一个开邻域 V，使得 \(\mu_{\mathrm{K}}(\mathrm{V} \cap \mathrm{K}) \leqslant \varepsilon\)；由 (i)，我们有 \(\mu(\mathrm{V} \cap \mathrm{K}) = \mu_{\mathrm{K}}(\mathrm{V} \cap \mathrm{K})\)，故 \(\mu(\mathrm{N}) = 0\)（由 \(\varepsilon\) 的任意性）。
\item
  对 K 中与 S 相交的每个开集 U，由假设有 \(\mu_{\mathrm{K}}(\mathrm{U}\cap\mathrm{S})\neq0\)，故由 (i) 得 \(\mu(\mathrm{U}\cap\mathrm{S})\neq0\)；又由于 \(U\cap S\) 在 S 中是开集，由 (i) 得 \(\mu_{\mathrm{S}}(\mathrm{U}\cap\mathrm{S})\neq0\)；由于 S 的每个非空开集都具有形状 \(U\cap S\)（其中 U 在 K 中为开集），这就证明了 \(\operatorname{Supp}(\mu_{\mathrm{S}})=\mathrm{S}\)。
\end{enumerate}

引理 3. --- 设 Y 是 X 的一个局部紧子空间；对 X 上的每个测度 \(\mu\)，有 \(|\mu_{\mathrm{Y}}| = |\mu|_{\mathrm{Y}}\) 。

设 \(f\) 是 \(\mathcal{K}_{+}(\mathrm{Y})\) 中的一个函数，\(\varepsilon\) 是任一数 \(>0\)；按定义，存在函数 \(g \in \mathcal{K}(\mathrm{Y};\mathbf{C})\)，使得 \(|g| \leqslant f\) 且 \(|\mu_{\mathrm{Y}}|(f) \leqslant |\mu_{\mathrm{Y}}(g)| + \varepsilon\) 。用 \(f'\) 与 \(g'\) 分别表示把 \(f\) 与 \(g\) 延拓成在 \(\mathrm{X} - \mathrm{Y}\) 上等于 0 所得的函数，我们有 \(\mu_{\mathrm{Y}}(g) = \mu(g')\)，且由于 \(|g'| \leqslant f'\) ，

\[
\left| \mu \left(g ^ {\prime}\right) \right| \leqslant | \mu | \left(\left| g ^ {\prime} \right|\right) \leqslant | \mu | \left(f ^ {\prime}\right) = | \mu | _ {Y} (f),
\]

由此得 \(|\mu_{\mathrm{Y}}|(f)\leqslant |\mu |_{\mathrm{Y}}(f) + \varepsilon\)，且由于 \(\varepsilon\) 任意，

\[
| \mu_ {\mathrm{Y}} | (f) \leqslant | \mu | _ {\mathrm{Y}} (f).
\]

另一方面，设 K 是 f 的支集，U 是 K 在 X 中的一个紧邻域，使得 \(|\mu|(\mathrm{U}-\mathrm{K}) \leqslant \varepsilon\) ；由 Urysohn 定理，存在函数 \(f_{1} \in \mathcal{K}_{+}(\mathrm{X})\)，它是 f 的延拓，其支包含于 U，且 \(\|f_{1}\| = \|f\|\) 。存在函数 \(h_{1} \in \mathcal{K}(\mathrm{X}; \mathbf{C})\)，使得 \(|h_{1}| \leqslant f_{1}\) 且 \(|\mu|(f_{1}) \leqslant |\mu(h_{1})| + \varepsilon\) 。若 h 是 \(h_{1}\) 在 Y 上的限制，则 \(h \in \mathcal{K}(\mathrm{Y}; \mathbf{C})\)，\(|h| \leqslant f\)，且 \(\mu(h_{1}) - \mu_{\mathrm{Y}}(h) = \mu(h_{1}\varphi_{\mathrm{U}-\mathrm{K}})\)，因此

\[
\left| \mu \left(h _ {1}\right) - \mu_ {\mathrm{Y}} (h) \right| \leqslant \| f \| \cdot | \mu | (\mathrm{U} - \mathrm{K}) \leqslant \varepsilon \| f \|;
\]

此外，同理有

\[
| \mu | (f _ {1}) - | \mu | _ {\mathrm{Y}} (f) = | \mu | (f _ {1} \varphi_ {\mathrm{U} - \mathrm{K}}) \quad \text { and } \quad \left| | \mu | (f _ {1}) - | \mu | _ {\mathrm{Y}} (f) \right| \leqslant \varepsilon \| f \|.
\]

由此我们推出

\[
| \mu | _ {\mathrm{Y}} (f) \leqslant | \mu_ {\mathrm{Y}} (h) | + \varepsilon (1 + 2 \| f \|) \leqslant | \mu_ {\mathrm{Y}} | (f) + \varepsilon (1 + 2 \| f \|)
\]

又由于 \(\varepsilon\) 任意，便得 \(|\mu|_{\mathrm{Y}}(f) \leqslant |\mu_{\mathrm{Y}}|(f)\)，证毕。

